\chapter{THEORY AND FUNDAMENTALS OF FLUORIDE ADSORPTION 
	USING GTMAC-MODIFIED COCONUT FIBER }

Fluoride adsorption using biomass-based adsorbents involves concepts from surface chemistry, environmental engineering, adsorption science, and statistical optimization. Understanding the fundamentals of adsorption mechanisms, surface modification techniques, adsorption isotherms, and experimental design methods is important for analyzing fluoride removal performance. This chapter discusses the theoretical background related to fluoride contamination, adsorption processes, GTMAC surface functionalization, batch adsorption studies, Langmuir isotherm modelling, and Design of Experiments (DOE). The chapter also explains the tools and techniques used for adsorption analysis and optimization in the present study. 

\section{Fundamentals of Fluoride Adsorption}

Fluoride contamination in drinking water is a major environmental issue affecting both rural and urban populations. Excess fluoride concentration above the WHO permissible limit of 1.5 mg/L can cause dental fluorosis and skeletal fluorosis. Adsorption is one of the most widely used techniques for fluoride removal because of its simplicity, low cost, and high efficiency. 

Adsorption is a surface phenomenon in which contaminants from a liquid phase accumulate on the surface of a solid material called an adsorbent. In the present work, GTMAC-modified coconut fiber acts as the adsorbent, while fluoride ions act as the adsorbate. The adsorption process mainly occurs through electrostatic attraction and ion-exchange mechanisms between fluoride ions and positively charged functional groups present on the modified adsorbent surface. 

\section{Surface Modification and GTMAC Functionalization}
Natural coconut fiber contains cellulose, hemicellulose, and lignin with hydroxyl functional groups capable of adsorption. However, the natural adsorption capacity of coconut fiber towards fluoride ions is limited due to low surface charge density and insufficient active adsorption sites. To improve adsorption efficiency, the coconut fiber is chemically modified using Glycidyl Trimethyl Ammonium Chloride (GTMAC). GTMAC introduces positively charged quaternary ammonium groups onto the adsorbent surface under alkaline conditions through etherification reactions. 

The surface modification process involves washing, drying, grinding, and sieving of coconut fiber followed by NaOH pretreatment for activation of cellulose hydroxyl groups. The activated material is then reacted with GTMAC at controlled pH and temperature conditions. The introduced quaternary ammonium groups enhance fluoride adsorption through electrostatic attraction and ion-exchange mechanisms. \gls{ftir} and \gls{sem} characterization techniques are used to confirm successful surface modification and surface morphology changes in the prepared adsorbent. 
 
\section{Batch Adsorption Studies}
Batch adsorption studies are performed to evaluate the fluoride removal efficiency of the prepared GTMAC-modified coconut fiber under different operating conditions. In batch adsorption experiments, a known amount of adsorbent is mixed with fluoride solution and agitated for a fixed period of time under controlled conditions. 

The major process parameters affecting adsorption include pH, adsorbent dose, initial fluoride concentration, and contact time. The residual fluoride concentration is measured using a Fluoride Ion-Selective Electrode F$^{-}$ ISE, and the adsorption performance is analyzed in terms of fluoride removal efficiency and adsorption capacity. Batch adsorption studies help in understanding adsorption equilibrium and the effect of different operating parameters on fluoride removal performance. 

\section{Design of Experiments (\gls{doe}) and Langmuir Modelling}
The Design of Experiments (\gls{doe}) is a statistical technique used to systematically plan and analyze experiments involving multiple variables. In the present study, Central Composite Design (CCD) under Response Surface Methodology (RSM) is used to optimize fluoride adsorption conditions by studying the interaction between pH, adsorbent dose, fluoride concentration, and contact time. 

The adsorption behaviour is further analyzed using the Langmuir adsorption isotherm model, which assumes monolayer adsorption on homogeneous adsorption sites. The Langmuir equation is expressed as: 
\begin{equation}
	q_e = \frac{q_{max} b C_e}{1 + b C_e}
	\label{eq:langmuir}
\end{equation}

where $q_e$ is the amount of fluoride adsorbed at equilibrium (mg/g), 
$q_{max}$ is the maximum adsorption capacity (mg/g), 
$b$ is the Langmuir constant related to adsorption energy (L/mg), 
and $C_e$ is the equilibrium fluoride concentration in solution (mg/L).

The Langmuir model is used to study adsorption characteristics and adsorption capacity of the modified coconut fiber under different operating conditions. DOE and Langmuir modelling together help in understanding adsorption behavior and optimizing fluoride removal efficiency. 

%\section{Use of Acronyms and Glossaries}
%% Acronyms
%\newacronym{doe}{DOE}{Design of Experiments}
%\newacronym{ccd}{CCD}{Central Composite Design}
%\newacronym{rsm}{RSM}{Response Surface Methodology}
%\newacronym{gtmac}{GTMAC}{Glycidyl Trimethyl Ammonium Chloride}
%\newacronym{ftir}{FTIR}{Fourier Transform Infrared Spectroscopy}
%\newacronym{sem}{SEM}{Scanning Electron Microscopy}
%\newacronym{fise}{F$^{-}$ ISE}{Fluoride Ion-Selective Electrode}
%\newacronym{ph}{pH}{Potential of Hydrogen}
%
%% Symbols
%\newglossaryentry{qe}{name=$q_e$, description={Adsorption capacity at equilibrium}, type=symbolList}
%\newglossaryentry{qmax}{name=$q_{max}$, description={Maximum adsorption capacity}, type=symbolList}
%\newglossaryentry{ce}{name=$C_e$, description={Equilibrium fluoride concentration}, type=symbolList}
%\newglossaryentry{c0}{name=$C_0$, description={Initial fluoride concentration}, type=symbolList}
%\vspace{0.75cm}
%
% \textbf{The chapters should not end with figures, instead bring the paragraph explaining about the figure at the end followed by a summary paragraph.}
%
%After elaborating the various sections of the chapter (From Chapter 2 onwards), a summary paragraph should be written discussing the highlights of that particular chapter. This summary paragraph should not be numbered separately. This paragraph should connect the present chapter to the next chapter.
